\synctex=1
\documentclass[a4paper,11pt,svgnames]{book}
%\usepackage{xcolor}
\usepackage[usenames,dvipsnames,table]{xcolor}

\usepackage[utf8]{inputenc}
\usepackage{mathtools}
\usepackage{thesis}
\usepackage{calc}
\usepackage{float}
\usepackage{acronym}
% \setlength{\marginparwidth}{2cm}
% \setlength{\voffset}{-0.25in}
% \setlength{\parskip}{0.1in}

\usepackage{todonotes}
\newcommand{\todoin}[1][]{\todo[inline, #1]}

%\includeonly{5-architecture}
%\includeonly{6-use-case}

%\usepackage{draftwatermark}
%\SetWatermarkScale{4}
\usepackage{subfig}
 \usepackage{listings}
  \usepackage{courier}
 \lstset{
 		 basicstyle=\footnotesize\ttfamily,
         numbers=none,               % Ort der Zeilennummern
         numberstyle=\tiny,          % Stil der Zeilennummern
         %stepnumber=2,               % Abstand zwischen den Zeilennummern
         numbersep=5pt,              % Abstand der Nummern zum Text
         tabsize=2,                  % Groesse von Tabs
         extendedchars=true,         %
         breaklines=true,            % Zeilen werden Umgebrochen
         keywordstyle=\color{red},
    	 frame=single,         
 %        keywordstyle=[1]\textbf,    % Stil der Keywords
 %        keywordstyle=[2]\textbf,    %
 %        keywordstyle=[3]\textbf,    %
 %        keywordstyle=[4]\textbf,   \sqrt{\sqrt{}} %
         stringstyle=\color{white}\ttfamily, % Farbe der String
         showspaces=false,           % Leerzeichen anzeigen ?
         showtabs=false,             % Tabs anzeigen ?
         belowcaptionskip=5pt, 
         xleftmargin=2em, 
         xrightmargin=0.8em,
         %backgroundcolor=\color{lightgray},
         showstringspaces=false      % Leerzeichen in Strings anzeigen ?        
 }
 \lstloadlanguages{% Check Dokumentation for further languages ...
         %[Visual]Basic
         %Pascal
         %C
         %C++
         %XML
         %HP
         Java
 }
 


\colorlet{punct}{red!60!black}
\definecolor{background}{HTML}{FFFFFF}
\definecolor{delim}{RGB}{20,105,176}
\colorlet{numb}{magenta!60!black}

\lstdefinelanguage{json}{
    basicstyle=\scriptsize\ttfamily,
    numbers=left,
    numberstyle=\scriptsize,
    stepnumber=1,
    numbersep=4pt,
    showstringspaces=false,
    breaklines=true,
    frame=lrtb,
    backgroundcolor=\color{background},
    literate=
     *{0}{{{\color{numb}0}}}{1}
      {1}{{{\color{numb}1}}}{1}
      {2}{{{\color{numb}2}}}{1}
      {3}{{{\color{numb}3}}}{1}
      {4}{{{\color{numb}4}}}{1}
      {5}{{{\color{numb}5}}}{1}
      {6}{{{\color{numb}6}}}{1}
      {7}{{{\color{numb}7}}}{1}
      {8}{{{\color{numb}8}}}{1}
      {9}{{{\color{numb}9}}}{1}
      {:}{{{\color{punct}{:}}}}{1}
      {,}{{{\color{punct}{,}}}}{1}
      {\{}{{{\color{delim}{\{}}}}{1}
      {\}}{{{\color{delim}{\}}}}}{1}
      {[}{{{\color{delim}{[}}}}{1}
      {]}{{{\color{delim}{]}}}}{1},
}
    %\DeclareCaptionFont{blue}{\color{blue}} 

  %\captionsetup[lstlisting]{singlelinecheck=false, labelfont={blue}, textfont={blue}}
  \usepackage{caption}
%\DeclareCaptionFont{white}{\color{white}}
%\DeclareCaptionFormat{listing}{\colorbox[HTML]{6DBAFF}{\parbox{\textwidth}{\hspace{5pt}#1#2#3}}}

\DeclareCaptionFont{white}{\color{white}}
\DeclareCaptionFormat{listing}{\colorbox{gray}{\parbox{\textwidth-20pt}{#1#2#3}}\vspace{0.01cm}}
\captionsetup[lstlisting]{format=listing,labelfont=white,textfont=white}

\captionsetup[lstlisting]{format=listing,labelfont=white,textfont=white, singlelinecheck=false, margin=0pt, font={bf,footnotesize}}


\usepackage{url}

%\usepackage[spanish]{babel}
%\usepackage[T1]{fontenc}
\usepackage{tabularx}
\usepackage{graphicx}
%\usepackage[usenames,dvipsnames,table]{xcolor}
\usepackage{verbatim}
\usepackage[section]{placeins}
\usepackage{listings}
%\bibliographystyle{unsrt}


\newcommand{\authorname}{NOMBRE DEL AUTOR}
\newcommand{\tfgtitle}{FULL TITLE OF THE THESIS}
\newcommand{\tfgtitlees}{TÍTULO COMPLETO DEL TFG}
\newcommand{\tutor}{TUTOR}
\newcommand{\supervisor}{PONENTE}
\newcommand{\fecha}{MES 20XX}

\usepackage[pdftex,
    pdfauthor={\authorname},
            pdftitle={\tfgtitle},
            pdfsubject={Master Final Project},
            pdfkeywords={GSI},
            pdfproducer={PDFTex},
            colorlinks=true,linkcolor=black,citecolor=black,urlcolor=black,hypertexnames=false]{hyperref}

%\usepackage{longtable}
\setcounter{secnumdepth}{3}

\usepackage{framed}

%lstlisting
\usepackage{listings}

\lstdefinelanguage{JavaScript}{
  keywords={typeof, new, true, false, catch, function, return, null, catch, switch, var, if, in, while, do, else, case, break},
  keywordstyle=\color{blue}\bfseries,
  ndkeywords={class, export, boolean, throw, implements, import, this},
  ndkeywordstyle=\color{darkgray}\bfseries,
  identifierstyle=\color{black},
  sensitive=false,
  comment=[l]{//},
  morecomment=[s]{/*}{*/},
  commentstyle=\color{purple}\ttfamily,
  stringstyle=\color{red}\ttfamily,
  morestring=[b]',
  morestring=[b]"
}

\definecolor{darkblue}{rgb}{0.0,0.0,0.6}

\lstdefinestyle{listXML}{language=XML, basicstyle=\ttfamily\diny, extendedchars=true,  belowcaptionskip=5pt,xleftmargin=0.4em, xrightmargin=0.3em, numbers=none, frame=single, breaklines=true, breakatwhitespace=true, breakindent=0pt, emph={}, emphstyle=\color{red}, basicstyle=\small\ttfamily, columns=fullflexible, showstringspaces=false, commentstyle=\color{gray}\upshape,
morestring=[b]",
morecomment=[s]{<?}{?>},
morecomment=[s][\color{orange}]{<!--}{-->},
keywordstyle=\color{cyan},
stringstyle=\color{black},
tagstyle=\color{darkblue},
morekeywords={xmlns,version,type}
}



\lstdefinestyle{mono}{
   framesep=8px,
   extendedchars=true,
   basicstyle=\ttfamily,
   showstringspaces=false,
   showspaces=false,
   tabsize=2,
   breaklines=true,
   showtabs=false,
   xleftmargin=8pt,
   xrightmargin=8pt
}


\lstdefinestyle{commands}{
   framesep=8px,
   extendedchars=true,
   basicstyle=\ttfamily,
   showstringspaces=false,
   showspaces=false,
   tabsize=2,
   breaklines=true,
   showtabs=false,
   xleftmargin=8pt,
   xrightmargin=8pt
}

\lstdefinestyle{consola}
   {basicstyle=\scriptsize\ttfamily,
    backgroundcolor=\color{white},
    frame=lrtb,
    numbers=none,
    xleftmargin=4pt,
    xrightmargin=4pt
   }

% Allows to change the color of chapter headers
\definecolor{chapterdetails}{HTML}{00a9e0}

% \usepackage[sf,bf]{titlesec}
% \titleformat{\chapter}[display]
%   {\normalfont\Large\sffamily\raggedleft}
%   {\vspace{5cm}\MakeUppercase{\chaptertitlename}%
%     \rlap{ \resizebox{!}{1.5cm}{\thechapter} \color{chapterdetails}\rule{5cm}{1.5cm}}}
%   {10pt}{\Huge}[{\color{chapterdetails}\titlerule[0.8mm] }]
% \titlespacing*{\chapter}{0pt}{30pt}{20pt}

\usepackage[sf,bf]{titlesec}
\titleformat{\chapter}[display]
  {\normalfont\Large\sffamily\raggedleft}
  {\vspace{5cm}\MakeUppercase{\chaptertitlename}%
    \rlap{ \resizebox{!}{1.5cm}{\thechapter} \color{chapterdetails}\rule{5cm}{1.5cm}}}
  {10pt}{\Huge}[{\color{chapterdetails}\titlerule[0.8mm] }]
\titlespacing*{\chapter}{0pt}{30pt}{20pt}

%\titleformat{\section}{\large\sffamily\bfseries}{\thesection}{1em}{}


\newenvironment{chapterintro}
{% This is the begin code
\large\it
}
{% This is the end code
}


% Tick symbols
\newcommand{\tickYes}{\checkmark}
\newcommand{\tickNo}{\hspace{1pt}\ding{55}}
% Fancy header
\usepackage{fancyhdr}
%Fancy chapter cover style


% Fancy box
\usepackage{fancybox} 
\setlength{\fboxrule}{1 pt} \setlength{\fboxsep}{10pt} \setlength{\shadowsize}{3pt}

%Sky color definition

%% Portada
\usepackage{eso-pic,graphicx}
\usepackage{chngpage}
% \usepackage{tikz}
% \usepackage[top=0cm, bottom=0cm, outer=0cm, inner=0cm]{geometry}


\begin{document}
\newcommand\litem[1]{\item{\bfseries #1 }}
\renewcommand{\arraystretch}{1.5} %Makes tables less crammed

\newcommand\headcell[1]{%
  \multicolumn{1}{|c|}{\cellcolor{DodgerBlue}\bfseries\sffamily\textcolor{white}{#1}}
}

\acrodef{TAS}[TAS]{Task Automation Service}

%Cuadros por tablas
%\renewcommand{\listtablename}{Tables Index}
%\renewcommand{\tablename}{Table} 

% \renewcommand{•}{•}*{\lstlistingname}{List of X}
	
\pagenumbering{Roman}

\pagestyle{empty}
% \tikz[remember picture,overlay] \node[opacity=1,inner sep=0pt] at (current page.center){\includegraphics[height=\paperheight]{img/portada.png}};

\vspace*{5.5cm}

\begin{center}

{\Large\rm \textbf{ MÁSTER UNIVERSITARIO EN\\
INGENIERÍA DE TELECOMUNICACIÓN\\}}

\vspace{1.0cm}

{\Large\rm \textbf{TRABAJO FIN DE MASTER}}

\vspace{2cm}

{\Large\rm\textbf{\tfgtitle}}

\vspace*{\fill}

{\Large\rm\textbf{NOMBRE DEL AUTOR}}

{\Large \textbf{20XX}}
\vspace{1.0cm}
\end{center}
\AddToShipoutPictureBG*{\includegraphics[width=\paperwidth,height=\paperheight]{img/portada.png}}

\cleardoublepage
\thispagestyle{empty}
\vspace*{3\baselineskip}
{\large{\bf TRABAJO DE FIN DE MASTER}}
\vspace{0.5cm}

\begin{rm}
\begin{tabular}{p{3cm}p{10cm}}
\textbf{Título:} & \tfgtitlees \\ 

\textbf{Título (inglés):} & \tfgtitle \\ 
\textbf{Autor:} & \authorname \\ 
\textbf{Tutor:} & \tutor \\ 
\textbf{Ponente:} & \supervisor \\ 
\textbf{Departamento:} & Departamento de Ingeniería de Sistemas Telemáticos \\ 
\end{tabular} \end{rm} \vspace{1cm}

{\large{\bf MIEMBROS DEL TRIBUNAL CALIFICADOR}} \vspace{0.5cm}

\begin{rm}
\begin{tabular}{p{3cm}p{10cm}}
\textbf{Presidente:} & -----\\
\textbf{Vocal:} & -----\\
\textbf{Secretario:} & -----\\
\textbf{Suplente:} & -----
\end{tabular}
\end{rm}
\vspace{1cm}

{\large{\bf FECHA DE LECTURA:}}
\vspace{1cm}

{\large{\bf CALIFICACIÓN:}}
\pagestyle{empty}
\cleardoublepage
\vspace*{\baselineskip}
\begin{center}
	{\LARGE\rm\textbf{UNIVERSIDAD POLITÉCNICA DE MADRID}\\
	\vspace{1.0cm}
	 ESCUELA TÉCNICA SUPERIOR DE\\ INGENIEROS DE TELECOMUNICACIÓN
	  }  \\

	 {\Large\rm Departamento de Ingeniería de Sistemas Telemáticos\\
	 Grupo de Sistemas Inteligentes  }  \\

\begin{figure}[!htbp]
	\centering
    \includegraphics[width=0.7\textwidth]{img/logo_etsit.jpg}

\end{figure}
	\vspace{1.0cm}
	{\LARGE\rm TRABAJO DE FIN DE MASTER\\
	\vspace{2.0cm}
    \tfgtitle
	 \vspace{0.5cm}}
	 
	 \vspace{1.0cm}
     \Large\rm\textbf{\author}\\
	 \vspace{1.0cm}
     \fecha
\end{center}  

%\cleardoublepage
%
%\begin{tabular}{p{10cm}p{4cm}}
%&\\
%&\\
%&\\
%&\\
%&\\
%&\\
%&\\
%&\\
%&\\
%&\emph{Write cool quote here}\\
%&\\
%\end{tabular}
\cleardoublepage
\phantomsection
\chapter*{Resumen}
\addcontentsline{toc}{chapter}{Resumen}

\vfill
\textbf{Palabras clave:} 
\cleardoublepage
\phantomsection
\chapter*{Abstract}
\addcontentsline{toc}{chapter}{Abstract}

\vfill
\textbf{Keywords: } 
%\cleardoublepage
\phantomsection
\addcontentsline{toc}{chapter}{Acknowledgement}

\begin{center}
\textbf{\large Acknowledgement}
\end{center}

Acknowledgement
\cleardoublepage
\phantomsection
\chapter*{Agradecimientos}
\addcontentsline{toc}{chapter}{Agradecimientos}

A Gauss

\cleardoublepage
\phantomsection
\addcontentsline{toc}{chapter}{Contents} % para que aparezca en el indice de 
\tableofcontents % indice de contenidos


\cleardoublepage
\phantomsection
\addcontentsline{toc}{chapter}{List of Figures} % para que aparezca en el indice de contenidos
\listoffigures % indice de figuras

%\cleardoublepage
%\phantomsection
%\addcontentsline{toc}{chapter}{List of Tables} % para que aparezca en el indice de contenidos
%\listoftables % indice de tablas

\cleardoublepage


%Header style
\pagestyle{fancy}
\fancyhf{}
\fancyhead[RO]{\sffamily \slshape \rightmark}
\fancyhead[LE]{\sffamily \slshape \leftmark}
%\renewcommand{\footrulewidth}{0.4pt} % grosor de la línea del pie
\fancyfoot[OR,EL]{\rmfamily \thepage} % texto derecha del pie
\pagenumbering{arabic}

\chapter{Introduction}

\textit{
This chapter is going to introduce the context of the project, including a brief overview of all the different parts that will be discussed in the project. It will also break down a series of objectives to be carried out during the realization of the project. Moreover, it will introduce the structure of the document with an overview of each chapter.}



\clearpage

\section{Context}

In this section we present the context of the project. The idea is to summarize the points that will be treated in the project.

\clearpage

\section{Project goals}

Describe the different goals of the project


\clearpage

\section{Structure of this document}

The remaining of this document is structured as follows:

\textbf{\textit{Chapter 2}} ... (Chapter 1 is the Introduction)
\chapter{State of Art}
\label{chap:state-of-art}
\textit{Introduction of the state of art}

\clearpage

\section{Sentiment Analysis}

\clearpage

\section{Neural Networks}


\chapter{Enabling Technolgies}
\label{chap:enabling_technologies}

\textit{This  chapter  offers  a  brief  review  of  the  main  technologies  that  have  made  possible this  project,  as  well  as  some of the related published works. }
\clearpage

\section{Machine Learning Technologies}

\subsection{TensorFlow}

\chapter{Architecture and Methodology}
\label{chap:architecture}
\textit{This chapter presents the methodology used in this work. It describes the overall architecture of the project, with the connections between the different components involved on the development of the project.}

\clearpage

\section{Methodology}

\clearpage

\section{Architecture}



\chapter{Case study}
\label{chap:case-study}

\textit{In this chapter we are going to describe a selected use case. This section will discuss the two use case: Sentiment and Emotion Analysis applying Deep Learning techniques, which will include the datasets used, the first steps with the neural networks and the use of more advance techniques.}

\clearpage

\section{Techniques}

\section{Datasets}

\section{Results}
...

\chapter{Conclusions}
\label{chap:conclusions}
\textit{This chapter will describe the achieved goals done by the master thesis following some the key points developed in the project.}

\clearpage
\section{Achieved Goals}

The achieved goals for this project are the following ones:

\begin{itemize}
    \item ...
\end{itemize}

\clearpage
\section{Conclusion}
\label{sec:conclusion}

The project has fulfilled the proposed objectives, but we must refer to each one of these objectives in to draw a general conclusion.




% apendices del proyecto
\appendix
\chapter{Example}

This is the description of the appendix.
\phantomsection
\nocite{*}
\addcontentsline{toc}{chapter}{Bibliography}
%ieeetr
\bibliographystyle{plain}
{
\small
\bibliography{biblio/ref}
}

%\cleardoublepage
\end{document}
